% ======================================================================
%  52_appendix_seminar.tex — 附录：讨论班手册（十二讲怎么开）
%  底本：KW 2006（hep-th/0604151），12 节，印刷 pp.2--222；本册主链 S1--S8
%  本文件只写教学安排；凡「讲次切分、复盘讲、验收」一律为 \pC
% ======================================================================

\section{讨论班手册：十二讲怎么开}
\label{sec:seminar}

\begin{keyidea}
\textbf{本附录要交付什么。}前面各章把一篇 231 页的论文重排成八步主链 S1--S8；本附录把它变成一门课：十二次、每次两小时、每次一步。目标不是「读完」，而是「讲出来、推出来、画出来」——结课时每人应能不看笔记重画 L0 的〈逻辑流〉，并对每个箭头说出一句话依据。凡讲次切分、时间与验收都是教学建议（\pC），不是原文结论。
\end{keyidea}

\subsection{开课方式：两小时、三种角色、读三遍}

\paragraph{时间表。}前 20 分钟预习检查：每人交一张卡片，正面写「预习里最不懂的一句话」，反面回应上一讲作业；主讲抽三张念出。中间 50 分钟主讲，只讲骨架与关键推导。随后 40 分钟白板推导，两人一组，第 30 分钟必须停手看提示。最后 10 分钟收尾：记录员宣读三条结论、一条未决，主讲指定下次的两位质询人。\pC

\paragraph{三种角色，每讲轮换。}\textbf{主讲}把「三个必须记住的箭头」写在白板角上；\textbf{质询}两人，主讲后 10 分钟内提至少三问，其中一问必须问「这一步依据原文哪一页、哪一公式」；\textbf{记录}一人写一页纪要（三条结论、一条未决、一张手绘图）进共享文档。十二讲内每人各当一次主讲、质询、记录。\pC

\paragraph{一本书读三遍。}第一遍（开学前）只读 L0（10\_layer0）：记住 S1--S8 与 M1--M5，遇公式跳过，只看编号与「一步步在干什么」。第二遍（随讲次）读 L1（20\_layer1\_phys 与 21\_layer1\_geom）的对应卡片：说清每讲的输入、关键构造、结论、出口。第三遍（期中后按兴趣）进 L2（30--33\_flesh\_*）：每人选一个节点做 20 分钟报告，必须含一处「本册重建、原文没写」的地方。\pC

\begin{warn}
讨论班不是论文导读：每讲产出是一条能画出来的箭头，不是「读完了第几页」。三遍的分工也不能混——先读编号，再读输入输出，最后读技术细节与可信度。\pC
\end{warn}

\subsection{十二讲的两种排法：标准版与精简版}

\paragraph{排法 A：十二讲标准版。}讲 \(k\) 就是原文 §\,\(k\)，学生自己就能预习，最省备课成本。缺点是不平衡：§7 只 13 页、§8 只 7 页，§10 有 36 页；讲 7 与讲 8 又是同一节点 S6 的两半。适合第一次开课。\pC

\begin{center}
\footnotesize
\begin{tabular}{@{}p{0.75cm}p{4.15cm}p{4.0cm}p{1.2cm}p{4.2cm}@{}}
\toprule
\tabhead{讲次} & \tabhead{原文（节与印刷页）} & \tabhead{本书对应} & \tabhead{主链} & \tabhead{一句话目标} \\
\midrule
1 & §1（pp.~2--5） & L0 前两节 & G0 & 说清全文要证什么 \\
2 & §2（pp.~6--15） & L1 上 §2.1--2.2 & S1 & 算出对偶群并写 S-对偶 \\
3 & §3（pp.~16--37） & L1 上 §3；L2 之一 & S2 & 推出 \(\Psi\)，说清 \(t=\pm\ii\) \\
4 & §4（pp.~38--55） & L1 上 §4；L2 之二 & S3 & 讲出三个复结构与纤维化 \\
5 & §5（pp.~56--77） & L1 上 §5；L2 之一、之二 & S4 & 陈述并验证 M2 \\
6 & §6（pp.~78--102） & L1 下 §6.1--6.4 & S5 & 定义 eigenbrane 并算例 \\
7 & §7（pp.~103--115） & L1 下 §7.1--7.2 & S6 & 说明为什么必须 twisted \\
8 & §8（pp.~116--122） & L1 下 §8.1--8.2 & S6 & 一行算出电特征膜 \\
9 & §9（pp.~123--147） & L1 下 §9；L2 之三 & S7 & 写出 Hecke 修改空间 \\
10 & §10（pp.~148--183） & L1 下 §10；L2 之三 & S7 & 把字典落到方程与 OPE \\
11 & §11（pp.~184--206） & L1 下 §11；L2 之四 & S8 & 造 c.c.\ 膜并完成翻译 \\
12 & §12（pp.~207--222） & L1 下 §12；L2 之四；L3 & S8 & 用边界条件验收全链 \\
\bottomrule
\end{tabular}
\end{center}

\paragraph{排法 B：十讲精简版。}把 §2 与 §3 合成一讲（对偶 \(+\) 扭化：都在准备参数，物理到几何的转译尚未开始），把 §9 与 §10 合成一讲（同一断言的两套语言：几何修改 vs.\ 方程与积分），§12 单独作验收讲；为保住十讲，§11 拆成两讲（§11.1--11.2 造膜、§11.3 扭化）。两次复盘讲正好落在讲 4 后与讲 8 后。\pC

\begin{center}
\footnotesize
\begin{tabular}{@{}p{0.75cm}p{4.15cm}p{4.0cm}p{1.2cm}p{4.2cm}@{}}
\toprule
\tabhead{讲次} & \tabhead{原文（节与印刷页）} & \tabhead{本书对应} & \tabhead{主链} & \tabhead{一句话目标} \\
\midrule
1 & §1（pp.~2--5） & L0 前两节 & G0 & 说清全文要证什么 \\
2 & §2\(+\)§3（pp.~6--37） & L1 上 §2、§3；L2 之一 & S1、S2 & 从对偶群走到 \(\Psi\) \\
3 & §4（pp.~38--55） & L1 上 §4；L2 之二 & S3 & 讲出三个复结构与纤维化 \\
4 & §5（pp.~56--77） & L1 上 §5；L2 之一、之二 & S4 & 陈述并验证 M2 \\
5 & §6（pp.~78--102） & L1 下 §6.1--6.4 & S5 & 定义 eigenbrane 并算例 \\
6 & §7\(+\)§8（pp.~103--122） & L1 下 §7、§8 & S6 & 从离散 flux 到电特征膜 \\
7 & §9\(+\)§10（pp.~123--183） & L1 下 §9、§10；L2 之三 & S7 & 't Hooft 就是 Hecke \\
8 & §11.1--11.2（pp.~184--201） & L1 下 §11.1--11.2 & S8 & 造 c.c.\ 膜、算开弦代数 \\
9 & §11.3（pp.~201--206） & L1 下 §11.3；L2 之四 & S8 & 扭化到任意 \(\Psi\) \\
10 & §12（pp.~207--222） & L1 下 §12；L2 之四；L3 & S8 & 用边界条件验收全链 \\
\bottomrule
\end{tabular}
\end{center}

\paragraph{两套怎么选，还能怎么并、怎么拆。}首次开课用排法 A（讲次＝原文节号，预习不用另编导引）；有经验后改排法 B，把省下的两讲换成复盘讲。想更细就拆 §4（超凯勒结构／Hitchin 纤维化各一讲）或拆 §10（Bogomolny 与冒泡／OPE 与扩展方程各一讲），得 14 讲；想更粗可再并 §7 与 §8，得 9 讲。两条不能动：§5 与 §11 不能并（M2 与 M5 是两次独立转译），§12 不能并进 §11（它是检查不是构造）。\pC

\subsection{每一讲的固定六栏}

每讲六栏：\textbf{预习}与\textbf{读后作业}给学生，其余四栏给讲者。五个问题必须课前发出，不许讲完才发。页码一律用印刷页，与正文的 \verb|\anchor{src p.N}| 约定一致。以下讲义、复盘与验收部分统一用小一号排，便于整页复印。\pC

{\small\linespread{0.95}\selectfont
% ----------------------------------------------------------------------
\subsection{第 1 讲：\S1（pp.~2--5）——总问题与六个主要想法}

\paragraph{预习。}\anchor{src p.2}读 pp.~2--5：p.~3 群对照表与 Atiyah 1977 提示，pp.~4--5 六个主要想法；配合 L0〈一页看懂〉。预先接受：\(\mathcal N=4\) 的存在性；几何 Langlands 陈述先当黑箱。\pC

\paragraph{本讲目标。}（1）能复述全文命题并指出两边对象类型；（2）能给六个想法各配一个 S 节点；（3）能指认哪些是物理输入、哪些是数学输出。\pC

\paragraph{五个必答问题。}（1）Table 1 里为什么 \(\mathrm{SO}(2n+1)\) 的对偶是 \(\mathrm{Sp}(n)\)，\(\mathrm{SU}(N)\) 的对偶还是自己？（2）为什么先紧化到曲线 \(C\) 上再谈「膜」？（3）六个想法哪个是种子，缺它其余不成立？（4）几何 Langlands 两边各是什么对象，「Hecke 特征」在两边各指什么？（5）哪些是猜想、哪些已证、哪些是物理启发？\pC

\paragraph{现场推导（30 分钟）。}画两列对照表：左列 4d 对象（\(\mathcal N=4\)、S-对偶、Wilson 与 't Hooft 线算子、沿 \(C\) 紧化），右列几何对象（\(\MH\)、A／B 模型、膜、D-模）；六个想法按原文顺序填入，标注每步是同构、范畴等价还是猜想。\emph{提示}：先填再排逻辑顺序。\emph{卡住时}看 L0〈六个想法与骨架编号〉。\pC

\paragraph{读后作业。}200 字内答「本讲核心障碍是什么」，必须点名一个具体对象或步骤，不许写「哲学意义」。\pC

\paragraph{卡点预警。}（1）第一讲就想读精确命题，而 \(\mathrm{Bun}_G\)、D-模、局部系统都还没定义；（2）把 §1 当导言：它其实是全链索引。\pC

% ----------------------------------------------------------------------
\subsection{第 2 讲：\S2（pp.~6--15）——\texorpdfstring{\(\mathcal N=4\)}{N=4} 与对偶群同构}

\paragraph{预习。}\anchor{src p.6}读 pp.~6--15：§2.1（pp.~6--11）场内容，§2.2（pp.~11--15）S-对偶与相位公式。预先接受：Montonen--Olive 对偶的陈述；\(SL(2,\ZZ)\) 对 \(\tau\) 的作用；\(\mathcal N=4\) 场内容。\pC

\paragraph{本讲目标。}（1）能写出含 \(n_g\) 的 S-对偶变换并说明它为何不是经典对称；（2）能从根与余根数据算出 \(\Ldual\)；（3）能写出中心荷与相位变换式 (2.22)--(2.25)。\pC

\paragraph{五个必答问题。}（1）\(S:\tau\mapsto-1/(n_g\tau)\) 里的 \(n_g\) 为何只在非单连通群出现且只取 \(2,3\)？（2）S-对偶同时交换哪三样东西（群、电荷、超荷相位）？漏掉哪样会让第 5 讲无法陈述镜像？（3）为什么磁单极子属于 \(\Ldual\) 的规范场？「GNO 对偶群与 Langlands 对偶群是同一个群」在哪一步用到？（4）(2.25) 的平方根为何说明 S 作用在超荷上是 \(SL(2,\ZZ)\) 的双覆盖？（5）为什么 S 只能作量子对称？\pC

\paragraph{现场推导（30 分钟）。}（a）从根系算 \(\Ldual\)：\(\Ldual\) 的根 \(=\) \(G\) 的余根，Cartan 矩阵转置；对 \(\mathrm{SU}(3)\)、\(\mathrm{Sp}(2)\) 验证 Table 1，再确认「对偶的对偶」回到 \(G\)。（b）取 \(\gamma=S\)，验证 (2.22) 的中心荷只差相位，并与 (2.25) 对上。\emph{提示}：Dynkin 图箭头反转，相位即 (2.25)。\emph{卡住时}看 L1 上 §2.2〈S-对偶：对偶群同构〉。\pC

\paragraph{读后作业。}算出 \(\mathrm{SO}(7)\) 的对偶群并验证二次对偶；200 字说明「S-对偶不是 \(\tau\mapsto-1/\tau\)」。\pC

\paragraph{卡点预警。}（1）把 \(\Ldual\) 当成「另一个群」：它是同一套根数据换一种读法，二次对偶才回到 \(G\)；（2）经典与量子对称的界限：\(T\) 保持经典拉氏量、\(S\) 不保持，这是 S1 的全部困难。\pC

% ----------------------------------------------------------------------
\subsection{第 3 讲：\S3（pp.~16--37）——扭化成族与典范参数 \texorpdfstring{\(\Psi\)}{Psi}}

\paragraph{预习。}\anchor{src p.16}读 pp.~16--37；时间紧只读 §3.1（pp.~17--24）与 §3.5（pp.~35--37），后者给出全文真正用的参数。预先接受：\(SU(4)_R\cong\mathrm{Spin}(6)\) 的表示论；三种扭化的定义；\(K=\Gamma_{89}\) 的本征值分配。\pC

\paragraph{本讲目标。}（1）能写出 \(\mathbb{CP}^1\) 族与参数 \(t\)；（2）能推出 \(\Psi\) 的公式并验证 \(t=\pm\ii\) 时 \(\Psi=\infty\) 且与 \(\tau\) 无关；（3）能说明消失定理为何让整族成为同一台拓扑机器。\pC

\paragraph{五个必答问题。}（1）为什么 \(t=\pm\ii\) 时 \(\Psi\) 与 \(\tau\) 无关？把 \(t^{2}=-1\) 代入 (3.50) 得什么？（2）两个不变超荷在 \(K=\Gamma_{89}\) 上本征值同为 \(-\ii\)；若不同，整族会怎样？（3）消失定理的前提是什么？为何度规只出现在一个 \(\delta_t\)-恰当项里？（4）\(t=\pm1\Rightarrow\Psi=\Re\tau\)、\(|t|=1\Rightarrow\Psi\in\RR\) 各说明什么？（5）为什么说 \(\Psi\) 比 \(\tau\) 更本质？\pC

\paragraph{现场推导（30 分钟）。}由 §3.4 拓扑项系数 \(-\frac{\Psi}{4\pi\ii}\int\Tr F\wedge F\) 与 \(\delta_t\) 的变换律推出 $\Psi=\frac{\tau+\bar\tau}{2}+\frac{\tau-\bar\tau}{2}\cdot\frac{t-t^{-1}}{t+t^{-1}}$，再反解 \(t^{2}=-(\Psi-\tau)/(\Psi-\bar\tau)\)，最后用 \(t\mapsto e^{2\ii\hat\varphi}t\) 验证 \(\Psi\mapsto-1/(n_g\Psi)\)。\emph{提示}：\((\tau\pm\bar\tau)/2\) 是 \(\tau\leftrightarrow\bar\tau\) 的对称与反对称部分。\emph{卡住时}看 L2 之一〈§3.5 canonical parameter〉。\pC

\paragraph{读后作业。}取 \(\tau=2\ii\)，算 \(t=0,\pm1,\pm\ii\) 处的 \(\Psi\)，解释 \(t=\pm\ii\) 为何永远给 \(\infty\)；200 字写「这一族为什么必须存在」。\pC

\paragraph{卡点预警。}（1）把 \(\Psi\) 当成 \(\tau\) 的函数：固定 \(\tau\) 时它跑遍 \(\mathbb{CP}^1\)；（2）把消失定理读成「解不变」：它说的是上同调类与 \(t\) 无关，代表元可以变。\pC

% ----------------------------------------------------------------------
\subsection{第 4 讲：\S4（pp.~38--55）——\texorpdfstring{\(\MH\)}{MH} 的超凯勒结构与 Hitchin 纤维化}

\paragraph{预习。}\anchor{src p.38}读 pp.~38--55：§4.1（pp.~40--49）超凯勒商、§4.2（pp.~49--52）复结构 \(I\) 与 \(J\)、§4.3（pp.~52--55）Hitchin 纤维化。预先接受：Hitchin 方程 (4.1)；Donaldson--Uhlenbeck--Yau 定理。\pC

\paragraph{本讲目标。}（1）能说明紧化为何把 4d 问题变成以 \(\MH\) 为靶空间的 sigma 模型；（2）能写出三个复结构上 \(\MH\) 的三种解读并验证 \(\Omega_I=\omega_J+\ii\omega_K\)；（3）能复述 Hitchin 纤维化与完全可积性。\pC

\paragraph{五个必答问题。}（1）取 \(M=\Sigma\times C\) 后为何低能理论是 \(C\) 上 Hitchin 模空间的 sigma 模型？「\(\Sigma\) 远大于 \(C\)」用在哪？（2）(4.1) 与复平坦性 \(F(A+\ii\varphi)=0\) 什么关系？实部虚部怎么分？（3）同一流形在 \(I\) 中是 Hitchin 对模空间、在 \(J\) 中是 \(\pi_1(C)\to G_{\CC}\) 表示模空间：两个全纯结构怎么共存？（4）\(\Omega_I\) 为何是 \(I\) 的全纯辛形式，其实部为何恰是 \(\omega_J\)？（5）Hitchin 纤维是复环面，维数多少？为何让 \(\MH\) 可积？\pC

\paragraph{现场推导（30 分钟）。}用 \(IJ=-JI=K\) 与 \(\omega_J(\cdot,\cdot)=g(J\cdot,\cdot)\) 证明 \(\Omega_I=\omega_J+\ii\omega_K\) 在 \(I\) 下是 \((2,0)\) 且 \(\dd\Omega_I=0\)；再算 \(\dim_{\RR}\MH=4\dim_{\CC}\MC\)，与 \(\varphi\in H^{0}(C,\mathrm{End}\,E\otimes K_C)\) 的自由度核对。\emph{提示}：先做 \(G=\mathrm{SU}(2)\) 的显式例子。\emph{卡住时}看 L2 之二〈超凯勒结构〉。\pC

\paragraph{读后作业。}对 \(G=\mathrm{SU}(2)\)、\(C=\mathbb{CP}^1\) 算 \(\dim\MC\)、\(\dim\MH\) 与纤维维数；200 字说明为什么只有 \(I\) 让 \(\MH\) 成为 Higgs 丛模空间。\pC

\paragraph{卡点预警。}（1）三个复结构的方向约定（\(I=\) Higgs 丛、\(J=\) 平坦联络、\(K=\) 中间）极易记反，而后文 A／B 型全依赖它；（2）\(\MH=T^*\MC\) 只在稳定点附近成立。\pC

% ----------------------------------------------------------------------
\subsection{第 5 讲：\S5（pp.~56--77）——扭化 TFT、镜像对称与 S-对偶}

\paragraph{预习。}\anchor{src p.56}读 pp.~56--77，重点是 §5.3（pp.~70--72）的 Table 2 与 \(\Psi=0,\infty\) 的模型识别；§5.1 的 \((J_+,J_-)\) 族、§5.2 广义复几何可略读。预先接受：广义复几何的基本定义；SYZ 猜想的陈述。\pC

\paragraph{本讲目标。}（1）能写出 \(\Psi=\infty\) 与 \(\Psi=0\) 分别对应哪个复结构上的 A／B 模型；（2）能陈述并验证 M2，说明为何需要复合一次超凯勒旋转 \(I\)；（3）能复述 SYZ 图像里「基上恒等、纤维上 T-对偶」。\pC

\paragraph{五个必答问题。}（1）\(\Psi=\infty\) 是 \(t=\ii\)，给 \(J\) 的 B 模型；施加 S 后得到的是 \(K\) 的 A 模型，为何不是直接得到 \(J\) 的 A 模型？（2）复合 \(I:\varphi\mapsto\ii\varphi\) 得 \(B(J)\leftrightarrow A(J)\)，原文为何说这一步不自然？（3）(5.35) 的 \(n_g\) 为何在非对角元上？单连通群时 S 才精确是 \(SL(2,\ZZ)\) 的 \(S\)？（4）Table 2 的六个变换里哪两个被 S 保持？（5）SYZ 图像中 S-对偶在 Hitchin 基上为何是恒等？\pC

\paragraph{现场推导（30 分钟）。}从 (5.35) 出发填 Table 2：先算 S 对 \(\Psi\) 的作用，再取 \(\Psi=\infty\) 与 \(\Psi=0\) 读出 \((w_+,w_-)\)，得 \(B(J)\leftrightarrow A(K)\) 与 \(A(J)\leftrightarrow B(K)\)，最后施加 \(I\) 得 \(B(J)\leftrightarrow A(J)\)。\emph{提示}：\(t=\ii\) 时 \((w_+,w_-)=(-\ii,-\ii)\) 给 \(J\) 的 B 模型，\(t=1\) 时 \((-1,1)\) 给 \(K\) 的 A 模型。\emph{卡住时}看 L1 上 §5.3〈特殊化：\(\Psi=0,\infty\) 与镜像对称〉。\pC

\paragraph{读后作业。}200 字说明 M2 依赖 \(\MH\) 的哪个结构；若只把 S-对偶当「两个 sigma 模型等价」会漏掉什么？\pC

\paragraph{卡点预警。}（1）把 S-对偶直接等同于镜像：少了一次 \(I\)，而 \(I\) 会把 \(A(J)\) 移出 \(\mathbb{CP}^1_g\) 族；（2）\(w_\pm\) 判据记混：\(I\) 上的 B 型与 \(J\) 上的 B 型不是同一件事。\pC

% ----------------------------------------------------------------------
\subsection{第 6 讲：\S6（pp.~78--102）——线算子等于膜算子，eigenbrane 的定义}

\paragraph{预习。}\anchor{src p.78}读 pp.~78--102：§6.1--6.2（pp.~79--89）拓扑 Wilson 与 't Hooft 算子、§6.3（pp.~89--92）紧化到二维、§6.4（pp.~93--102）为本讲核心。预先接受：拓扑算子定义；沿 \(C\) 紧化得二维 sigma 模型；膜范畴及其模结构。\pC

\paragraph{本讲目标。}（1）能写出 eigenbrane 定义 \(XB\simeq B\otimes V\) 并解释两边；（2）能说明「来自 4d」给出的两个自动结构；（3）能指出 Wilson 族与 't Hooft 族的 OPE 各对应哪个表示环。\pC

\paragraph{五个必答问题。}（1）\(XB\simeq B\otimes V\) 中的 \(V\) 是「另一个膜」还是「一个向量空间」？为何必须张量积而不能直和？（2）为何来自 4d 的线算子自动交换？脚注 31 在哪一步用到「支撑点可在 \(C\) 中分开」？（3）特征值为何构成 \(C\) 上平坦丛？「移动 \(p\) 只差 BRST 平凡项」用在哪？（4）Wilson 族与 't Hooft 族的 OPE 各是什么代数？后者原文的把握属哪一档？（5）special Wilson 与 special 't Hooft 的类型 \((B,B,B)\)、\((B,A,A)\) 怎么定出？\pC

\paragraph{现场推导（30 分钟）。}取零维膜 \(B_x\)（\(x\in\MH\)），把拓扑 Wilson 线算子贴着边界插入，算出 \(XB_x=E_q(R)|_x\otimes B_x\)，读出 (6.22) 中 \(V\) 的谱与维数；再用 S-对偶说明 magnetic eigenbrane 属 't Hooft 族。\emph{提示}：Wilson 线只带电荷 \(q\)，「作用」就是在 \(x\) 处取值。\emph{卡住时}看 L1 下 §6.4〈边界附近的线算子与 eigenbrane〉。\pC

\paragraph{读后作业。}200 字说明为何 \(V\) 不能换成直和；补小计算：取 \(G=U(1)\)、电荷 \(q\) 的表示，写出 \(V\)。\pC

\paragraph{卡点预警。}（1）把 \(V\) 当成膜：它是纤维／向量空间；（2）(6.22) 是定义不是定理：真正内容是「来自 4d」给两个自动结构。\pC

% ----------------------------------------------------------------------
\subsection{第 7 讲：\S7（pp.~103--115）——离散 flux 与 universal bundle 的障碍}

\paragraph{预习。}\anchor{src p.103}读 pp.~103--115：§7.1（pp.~103--109）中心完美配对与 universal bundle 障碍、§7.2（pp.~110--115）flux 的二维含义。预先接受：\(Z=\Lambda^\vee/\Lambda\) 与中心扩张；twisted 向量丛定义；Künneth 分解。\pC

\paragraph{本讲目标。}（1）能说明 universal \(G\)-bundle 为何不存在；（2）能写出 twist 标签与障碍类的关系；（3）能复述「正确对象是所有 \(e_0\) 的 twisted zerobrane 之和」。\pC

\paragraph{五个必答问题。}（1）为何伴随群 \(G_{\mathrm{ad}}\) 有 universal bundle 而单连通 \(G\) 没有？障碍出现在三重重叠的哪一步？（2）twist 的标签是中心元素，它与离散电 flux \(e_0\) 是什么关系？为何是同一个东西的两种说法？（3）配对 \(\Upsilon(a,b)=\exp(2\pi\ii\langle a,b\rangle)\) 为何使 \(H^{2}(M,\ZZ)\) 自对偶？（4）S-对偶在离散 flux 上为何是 \((e,m)\mapsto(-m,e)\)？（5）若跳过本讲直接读 §8.2，会漏掉什么？\pC

\paragraph{现场推导（30 分钟）。}写出 universal bundle 拼接的障碍：过渡函数在三重重叠上只差中心元素 \(f_{ijk}\in Z\)，给出障碍类 \(\zeta=\xi(E_{\mathrm{ad}})\in H^{2}(M,\ZZ)\)；再说明 twisted bundle 用 \(\varrho(\zeta)\in H^{2}(M,U(1))\) 扭一下就能存在。\emph{提示}：把不同平凡化的差别看成 1-上链，再算其上边缘。\emph{卡住时}看 L1 下 §7.1。\pC

\paragraph{读后作业。}取 \(G=\mathrm{SU}(2)\)，说明 \(e_0\) 的两个取值给出两个 twisted zerobrane，并解释为何必须求和；200 字回答「twist 是拓扑还是背景场」。\pC

\paragraph{卡点预警。}（1）把 twisted 当成「加背景磁通」：它是丛在模空间上的拓扑，不是新的场；（2）把「求和」当技术细节：它才是 S-对偶不变性的来源。\pC

% ----------------------------------------------------------------------
\subsection{第 8 讲：\S8（pp.~116--122）——zerobrane 就是电特征膜}

\paragraph{预习。}\anchor{src p.116}读 pp.~116--122：§8.1（pp.~116--120）Wilson 算子如何作用于膜、§8.2（pp.~121--122）为本讲核心。预先接受：Wilson 算子的作用规则；\(\MH\) 上点的几何；上一讲的 twisted bundle。\pC

\paragraph{本讲目标。}（1）能写出 \(E_p(R)\) 的定义与它的扭化；（2）能用「任意向量丛在一点处平凡」一行算出 zerobrane 是电特征膜；（3）能说明为何「所有 \(e_0\) 求和」给出 S-对偶不变的对象。\pC

\paragraph{五个必答问题。}（1）§8.2 的一行计算为何合法？它忽略了什么、这些忽略在哪一步补回？（2）位于 \(x\in\MH\) 的 zerobrane 的特征值为何是 \(E_q(R)|_x\)？\(R\) 与 \(q\) 各是什么？（3）Wilson 线把离散电场改成 \(e_0\to e_0+\theta_R\)，这一步在 (8.1)--(8.3) 里怎么体现？（4）为何电特征膜是一个求和而非单个 zerobrane？（5）若没有第 7 讲的 twist 修正，§8.2 会错成什么样？\pC

\paragraph{现场推导（30 分钟）。}取 \(B_x\) 为 \(x\) 处 zerobrane，把 Wilson 算子沿 \(C\) 中圈插入，用「\(E\) 在 \(x\) 处可平凡化」把 \(W_q(R)|_x\) 化成 \(E_q(R)|_x\)，再验证 \(E_p(R)\) 随 \(p\) 的平坦性。\emph{提示}：只有表示 \(R\) 的权（电荷 \(q\)）出现。\emph{卡住时}看 L1 下 §8.2。\pC

\paragraph{读后作业。}200 字说明「§8.2 证明了什么、没证明什么」，各写一条。\pC

\paragraph{卡点预警。}（1）以为「zerobrane \(=\) 电特征膜」显然：它依赖第 7 讲的求和修正；（2）把 \(E_q(R)\) 当数值而不是丛。\pC

% ----------------------------------------------------------------------
\subsection{第 9 讲：\S9（pp.~123--147）——\texorpdfstring{'t Hooft}{'t Hooft} 算子等于 Hecke 算子}

\paragraph{预习。}\anchor{src p.123}读 pp.~123--147：§9.1（pp.~128--134）静态 't Hooft 与 Hecke 修改、§9.2（pp.~135--144）修改空间、§9.3（pp.~145--147）仿射 Grassmann。预先接受：Hecke 修改的代数几何定义 \((E',\beta)\)；\(\mathrm{Gr}_G=G((z))/G[[z]]\) 的胞腔分解。\pC

\paragraph{本讲目标。}（1）能写出 't Hooft 算子标签 \(\rho:U(1)\to G\) 与 GNO 荷的对应；（2）能写出修改的双陪集分类 \(G(\mathcal O)\backslash G(\mathcal K)/G(\mathcal O)\)；（3）能说明修改空间与 Schubert 胞腔的关系。\pC

\paragraph{五个必答问题。}（1）'t Hooft 算子是 4d 对象、Hecke 修改是 \(C\) 上一点的代数操作，「等于」在哪一步被定义出来？（2）为何修改数据是 \((E',\beta)\) 而不只是 \(E'\)？\(\beta\) 的失效程度怎么给出离散类型？（3）双陪集为何按支配余权分类？连续参数又是什么？（4）\(GL_2\) 的型 \((1,0)\) 修改空间是 \(\mathbb{CP}^1\)，型 \((2,0)\) 的紧化为何自动包含型 \((1,1)\)？（5）'t Hooft 算子的路径积分为何等于在修改空间上积分？\pC

\paragraph{现场推导（30 分钟）。}算 \(GL_2\)、型 \((1,0)\) 的修改空间：\(E'\) 由 \(z^{-1}s_1,s_2\) 生成，截面 \(s(z)=cz^{-1}s_1(0)+(h_1,h_2)\)，参数为 \(s_1(0)\in\CC^{2}\setminus\{0\}\) 模 \(\CC^{*}\)，故空间 \(\cong\mathbb{CP}^1=\mathrm{Gr}(1,2)\)；再算型 \((2,0)\) 的两级极点，得 \(\mathbb{WP}^{2}(1,1,2)\) 挖一点。\emph{提示}：先写局部标架再写截面。\emph{卡住时}看 L2 之三〈第二层：Hecke 修改的数学定义〉。\pC

\paragraph{读后作业。}对 \(GL_2\) 写出型 \((1,1)\) 的修改并验证它等于 \(E\otimes\mathcal O(x)\)；200 字写「't Hooft \(=\) Hecke 的可检验内容」，必须给一个可算的量。\pC

\paragraph{卡点预警。}（1）把「等于」读成严格同构：它是一条字典，可信度按 L0〈可信度地图〉分档；（2）把「局部修改」读成「全局不变」：\(c_1\) 可变，稳定性也可能变。\pC

% ----------------------------------------------------------------------
\subsection{第 10 讲：\S10（pp.~148--183）——Bogomolny 方程、冒泡与 OPE}

\paragraph{预习。}\anchor{src p.148}读 pp.~148--183：§10.1（p.~148）边界条件与指数定理、§10.2（p.~155）单极子成泡、§10.3（p.~159）Kähler 商、§10.4（p.~167）OPE、§10.5（p.~179）扩展方程。预先接受：Bogomolny 方程与单极子模空间；Kähler 商构造。\pC

\paragraph{本讲目标。}（1）能写出 't Hooft 算子路径积分的边界条件（一端 Dirichlet、一端 Neumann）；（2）能说明单极子冒泡 \(=\) 修改空间非紧的边界；（3）能复述 OPE 系数等于表示论重数、扩展方程只有 \(E\) 被修改。\pC

\paragraph{五个必答问题。}（1）为何 't Hooft 算子的路径积分是「在非紧空间上积分」？非紧性从哪来？（2）单极子冒泡为何正好给出修改空间的边界？冒泡掉的单极子去了哪里？（3）Kähler 商等于复商这一步，为何使修改空间出现在 Bogomolny 模空间中？（4）(10.4) 的 OPE 系数等于表示论重数：两边各由什么决定？它给第 9 讲哪句断言一个「表示论版本」？（5）扩展方程里只有 \(E\) 被修改、\(\varphi\) 不能，几何含义是什么？\pC

\paragraph{现场推导（30 分钟）。}写出 't Hooft 算子附近的 Bogomolny 方程与边界条件（\(x^{1}=0\) 处一端 Dirichlet、一端 Neumann），说明解在远处趋于真空、在奇点处趋于单极子；再写扩展情形下 \(\varphi\) 的边界条件。\emph{提示}：先做 \(G=\mathrm{SU}(2)\) 的单极子。\emph{卡住时}看 L2 之三〈第五层：Bogomolny 方程〉。\pC

\paragraph{读后作业。}200 字说明「为何 §10 要把 't Hooft \(=\) Hecke 再证一遍」；另算型 \(\lambda\) 的修改空间复维数公式。\pC

\paragraph{卡点预警。}（1）§10 长但结构单一：它是把 §9 的字典落到方程上，容易读成一堆独立计算；（2）冒泡不是「丢掉单极子」，而是把不稳定边界收进修改空间。\pC

% ----------------------------------------------------------------------
\subsection{第 11 讲：\S11（pp.~184--206）——典范余各向同性膜与膜到 D-模的翻译}

\paragraph{预习。}\anchor{src p.184}读 pp.~184--206：§11.1（pp.~185--196）c.c.\ 膜与 (11.1)、§11.2（pp.~196--201）开弦代数 \(=\) D-模、§11.3（pp.~201--206）扭化。预先接受：开闭拓扑场论公理；D-模与自旋丛 \(K^{1/2}\)；B-field 的规范变换。\pC

\paragraph{本讲目标。}（1）能说明翻译器为何必须空间充满；（2）能写出 (11.1) 并说明它给的是可积性而不是膜的存亡；（3）能说出 \(\Psi=0\) 与 \(\Psi\neq0\) 的代数层分别是 \(\Dmod_{K_M^{1/2}}\) 与 \(\Dmod_{K_M^{1/2}\otimes L^{\Psi}}\)。\pC

\paragraph{五个必答问题。}（1）为何普通 A-膜（Lagrangian 支撑）不能当翻译器？「空间充满」是必要还是充分条件？（2）(11.1) 的功能是让膜存在，还是让翻译可积？两者区别在哪？（3）为何开弦代数必须逐点同构于 \(\Dmod_{K_M^{1/2}}\)，而不能只是某个 \(\Dmod_L\)？时间反演用在哪一步？（4）\(F=(\Im\tau)\omega_J\) 是平凡线丛上的固定联络，「拓扑平凡」为何反而关键？（5）非整数 \(\Psi\) 时线丛 \(L^{\Psi}\) 不存在，扭化 D-模层怎么定义？\pC

\paragraph{现场推导（30 分钟）。}取 \(X=\MH\)、\(K\) 中的 A-模型：解 (11.1) 的 \(F\) 并验证 \(\hat N=\omega^{-1}F=I\)；再说明 c.c.\ 膜的开弦代数为何给出 \(\Dmod_{K_M^{1/2}}\)（用胶合修正与 \(\Omega_I\) 的 Poisson 括号）。\emph{提示}：先算 \(g+F\) 的反对称部分。\emph{卡住时}看 L2 之四〈典范余各向同性 A-膜〉。\pC

\paragraph{读后作业。}200 字说明「为何说翻译是识别而不是量子化」；补一例：取 \(\varphi=0\) 的膜，验证 \((B_{c.c.},B')\) 弦的层就是 \(\Dmod_{K_M^{1/2}}\) 的定义模。\pC

\paragraph{卡点预警。}（1）把「空间充满」当成「边界条件自由」：恰恰相反，c.c.\ 膜是混合边界条件；（2）以为存在一个叫「量子化」的函子：原文走的是「算 endomorphism 代数」这条路。\pC

% ----------------------------------------------------------------------
\subsection{第 12 讲：\S12（pp.~207--222）——三类膜与四维边界条件（验收讲）}

\paragraph{预习。}\anchor{src p.207}读 pp.~207--222：§12.1（pp.~208--210）一般方法、§12.2（pp.~211--214）\((B,B,B)\) 膜、§12.3（pp.~215--217）\((B,A,A)\) 膜、§12.4（pp.~217--221）\((A,B,A)\) 膜与 c.c.\ 膜。预先接受：超对称边界条件的标准分析（玻色边界项、费米边界条件、\(T\) 矩阵）；未破缺超对称判据。\pC

\paragraph{本讲目标。}（1）能列出三类膜及其 \((A,B,A)\) 类型；（2）能用边界条件认出第 11 讲的 c.c.\ 膜；（3）能说明 §12 是自洽性检查而非新构造，并说清为何 D-模必须住在 \(\mathrm{Bun}_G\) 上。\pC

\paragraph{五个必答问题。}（1）为何 \((B,B,B)\) 膜的两种四维实现只在奇点附近才不同？「提升不唯一」是缺陷还是特征？（2）S-对偶如何把 \((B,B,B)\) 变成 \((B,A,A)\)？为何原文只变换未破缺超对称、不追具体膜？（3）\((A,B,A)\) 膜的边界条件里哪个 \(T\) 矩阵起作用？它为何能认出 c.c.\ 膜？（4）什么迫使 D-模住在 \(\mathrm{Bun}_G\) 而不是 \(\MC\) 上？（5）§12 的结论是「新定理」还是「检查」？若为检查，它检查了什么？\pC

\paragraph{现场推导（30 分钟）。}取 \(T=\Gamma^{04}\)，由费米边界条件 \(T\lambda|=\lambda|\) 与 \((T-1)\delta\lambda=0\)、\(\delta\lambda=\frac12\Gamma^{IJ}F_{IJ}\epsilon\)，推出 \([T,\Gamma^{IJ}F_{IJ}]\epsilon=0\)，即自对偶条件；再对 \((A,B,A)\) 写出混合边界条件。\emph{提示}：先按 \([T,\Gamma^{IJ}]\) 是否为零分类。\emph{卡住时}看 L2 之四〈第 12 节：检查清单〉。\pC

\paragraph{读后作业。}把三类膜按 \((I,J,K)\) 画成 \(3\times3\) 表并标出 S-对偶作用；200 字说明「为何 §12 被称为诚实度指标」。\pC

\paragraph{卡点预警。}（1）把类型三元组当成三种不同的膜理论：它们只是同一膜相对三个复结构的类型记录；（2）§12 长而结论少，产出只是「第 11 节用的膜确有四维来源」。\pC

\subsection{三次复盘讲：插在哪里、复盘什么}

\paragraph{复盘一：物理到几何（S1--S4 讲完）。}12 讲版插在第 5 讲后，10 讲版插在第 4 讲后。90 分钟：分组重画从 4d 到 2d 的图，每个箭头须注「依据哪一页、哪一公式」；交叉质询，每组解释两个箭头。自测两问：为什么需要 \(\Psi\)？\(\Psi=\infty\) 的 \(B(J)\) 与 \(\Psi=0\) 的 \(A(K)\) 分属哪个群？\pC

\paragraph{复盘二：算子到膜到 D-模（S5--S8 的转译讲完）。}12 讲版插在第 11 讲后，10 讲版插在第 8 讲后（翻译已完成，只剩 §11.3 的推广与 §12 的检查）。画 Wilson／'t Hooft 算子 \(\to\) eigenbrane \(\to\) twisted zerobrane／Hitchin 纤维 \(\to\) c.c.\ 膜 \(\to\) D-模，并区分三个「等」的强度：第 6 讲是定义，第 9--10 讲是字典，第 11 讲是定理级识别。反例环节：每人带一个「自以为懂了其实错了」的地方，用 L0〈概念缠绕〉核对。\pC

\paragraph{复盘三：全链与 2024 年证明的情况。}前半把 S1--S8 与 M1--M5 逐条对账，每条注明本书哪一节、原文哪一页、可信度哪一档。后半讲 2024 年以来 Gaitsgory--Raskin 等五篇系列预印本：证明特征 \(0\) 代数闭域上光滑射影曲线的几何 Langlands（D-模、未分歧情形）。要讲清三点：这是数学定理，不是本册物理路线的证明；数学走范畴化路线，与物理解释不互相验证；§12 的检查清单与定理适用范围恰好一致。\pC

\subsection{验收标准与十分钟口试}

\paragraph{通过标准。}能不看笔记重画 L0〈逻辑流〉，对每个箭头说出一句话依据（哪一步、哪一页、哪一公式），并指出至少两个箭头属「本册重建」（\pB）而非「原文明确」（\pA）；可以卡住一个箭头，但要知道卡在哪。\pC

\paragraph{十分钟口试清单（六问）。}（1）一句话说出全文主命题与两边对象类型。（2）为何进入几何 Langlands 陈述的是 \(\Psi\) 而不是 \(\tau\)？\(\Psi=0,\infty\) 各对应什么？（3）eigenbrane 的定义是什么？特征值为何构成平坦丛？（4）'t Hooft \(=\) Hecke 的可检验内容是什么？给一个可算的量。（5）c.c.\ 膜为何必须空间充满？(11.1) 起什么作用？（6）§12 为何是「诚实度指标」？举一个它检查出的问题。\pC

\paragraph{评分与补考。}六问中四问完整、其余能说出卡点即通过；否则重读 L0 对应小节，两周内再约。\pC

\begin{warn}
验收不要放水成「能背出 S1--S8 的名字」：要考的是箭头——为何这一步接上下一步、依据在哪一页、哪句原文哪句我们补的。把 \pA 与 \pB 混着说的答案按未通过处理（对讲者同样适用）。\pC
\end{warn}
}% end \small
